% !TEX program = lualatex
\documentclass[professionalfonts]{beamer}
\newcommand{\slidemode}{1}  % カラー投影モード
\usepackage{beamer_template}
\usepackage{enumerate}

\newcommand{\LectureNum}{10}
\newcommand{\LectureTitle}{一般の周期関数のフーリエ級数（２）}
\newcommand{\CourseName}{応用数学}
\renewcommand{\TermName}{前期}

\begin{document}

\begin{frame}{第10回　一般の周期関数のフーリエ級数（２）}
  \begin{block}{本回の内容}
  一般の周期関数のフーリエ級数（２）
  \end{block}
  \vfill\hfill{\scriptsize \textcolor{gray}{第3章 フーリエ解析　§1}}
\end{frame}

\begin{frame}{定理：フーリエ余弦展開（偶関数）}
  \begin{block}{}
  $f(t)$ が周期 $2l$ の偶関数のとき
  \[
    f(t) = c_{0} + \sum_{n=1}^\infty a_{n}\cos(n\pi t/l)
  \]
  \[
    c_{0} = (1/l)\int_{0}^{l} f(t) dt,\quad a_{n} = (2/l)\int_{0}^{l} f(t)\cos(n\pi t/l) dt
  \]
  \end{block}
  {\footnotesize \textcolor{gray}{偶関数なので $b_{n}=0$、$[0,l]$ 上の積分のみで決まる}}
\end{frame}

\begin{frame}{定理：フーリエサイン展開（奇関数）}
  \begin{block}{}
  $f(t)$ が周期 $2l$ の奇関数のとき
  \[
    f(t) = \sum_{n=1}^\infty b_{n}\sin(n\pi t/l)
  \]
  \[
    b_{n} = (2/l)\int_{0}^{l} f(t)\sin(n\pi t/l) dt
  \]
  \end{block}
  {\footnotesize \textcolor{gray}{奇関数なので $c_{0}=0, a_{n}=0$、$[0,l]$ 上の積分のみで決まる}}
\end{frame}

\begin{frame}{例題4　（p.\,84）}
  \begin{exampleblock}{問題}
    $f(x) = x(1-x) (0 \leq x \leq 1), f(-x) = -f(x), f(x+2) = f(x)$ のフーリエ展開を求めよ\\[2pt]
  \end{exampleblock}
\end{frame}

\begin{frame}{例題4　解答}
  $f(x)$ は奇関数なので $c_{0} = 0, a_{n} = 0,$ フーリエ正弦展開を求めればよい\\[4pt]
  $b_{n} = 2\int_{0}^{1} x(1-x) \sin(n\pi x) dx = 4(1-(-1)^n)/(n^{3}\pi^{3})$
  \medskip\textbf{答}：$f(x) = (8/\pi^{3})(\sin \pi x + \sin(3\pi x)/27 + \sin(5\pi x)/125 + \ldots) = (8/\pi^{3})\sum_{k=0}^\infty \sin((2k+1)\pi x)/(2k+1)^{3}$
\end{frame}

\begin{frame}{例題5　（p.\,84--86）}
  \begin{exampleblock}{問題}
    次の等式が成り立つことを証明せよ $: 1 - 1/3 + 1/5 - \ldots = \pi/4$\\[2pt]
  \end{exampleblock}
\end{frame}

\begin{frame}{例題5　解答}
  第8回 例題2 の関数（周期2の矩形波）$f(x)=1\,(0\leq x<1), 0\,(-1\leq x<0)$ は $x=1/2$ で連続だから\\[4pt]
  $f(1/2) = 1/2 + (2/\pi)(\sin(\pi/2) + \sin(3\pi/2)/3 + \sin(5\pi/2)/5 + \ldots)$\\[4pt]
  $= 1/2 + (2/\pi)(1 - 1/3 + 1/5 - \ldots)$\\[4pt]
  $f(1/2) = 1$ なので $: 1 = 1/2 + (2/\pi)(1 - 1/3 + 1/5 - \ldots)$
  \medskip\textbf{答}：$1 - 1/3 + 1/5 - \ldots = \pi/4$
\end{frame}

\begin{frame}{演習問題}
  \begin{exampleblock}{自分で解いてみよう}
  \begin{description}
    \item[問4] 次の関数のフーリエ級数を求めよ
    $(1) f(x) = \begin{cases} -1 & (-1 \leq x < 0) \\ 1 & (0 \leq x < 1) \end{cases},  f(x+2) = f(x)$\\
    $(2) g(x) = 1 - |x|  (-1 \leq x < 1),  g(x+2) = g(x)$\\
    $(3) h(x) = \cos(\pi x/4) \ (-2 \leq x < 2),  h(x+4) = h(x)$\\
    \item[問5] 次の問いに答えよ
    (1) 例題3の関数のフーリエ級数を用いて、次の公式を導け\\
    $1/1^{2} + 1/3^{2} + 1/5^{2} + \ldots = \pi^{2}/8$\\
    (2) 例題4の関数のフーリエ級数を用いて、次の公式を導け\\
    $1/1^{3} - 1/3^{3} + 1/5^{3} - \ldots = \pi^{3}/32$\\
  \end{description}
  \end{exampleblock}
\end{frame}

\end{document}
